% Include in main.tex with: % Include in main.tex with: % Include in main.tex with: % Include in main.tex with: \input{figures/fbar_strength_comparison}
% Requires in preamble: \usepackage{pgfplots}, \pgfplotsset{compat=1.18}
%                       \usepgfplotslibrary{groupplots}
\begin{figure}[htbp]
\centering
\begin{tikzpicture}
\begin{groupplot}[
    group style={
        group size=2 by 1,
        horizontal sep=1.8cm,
    },
    width=0.42\textwidth,
    height=0.42\textwidth,
    grid=major,
    grid style={gray!25},
    tick align=outside,
    xlabel={Max F-bar pressure correction (ref)},
    nodes near coords,
    nodes near coords style={font=\small, anchor=south west},
    point meta=explicit symbolic,
]

%% --- Left panel: new method ---
\nextgroupplot[
    title={New method},
    ylabel={Relative error at target},
    ymax=0.7
]
\addplot[only marks, mark=*, mark size=3pt, blue] table[
    col sep=comma,
    x=fbar_pressure_correction_max_ref,
    y=rel_err_target,
    meta=nu,
]{data/strength/fbar_strength_results_new.csv};

%% --- Right panel: old method ---
\nextgroupplot[
    title={Old method},
    ylabel={Relative error at target},
    yticklabel pos=right,
    ymax=0.7
]
\addplot[only marks, mark=*, mark size=3pt, red] table[
    col sep=comma,
    x=fbar_pressure_correction_max_ref,
    y=rel_err_target,
    meta=nu,
]{data/strength/fbar_strength_results_old.csv};

\end{groupplot}
\end{tikzpicture}
\caption{Dynamic simulation - Relative error at target vs.\ maximum F-bar pressure correction for the new method (left) and old method (right), with Poisson's ratio $\nu$ labelled at each point.}
\label{fig:fbar_strength_comparison}
\end{figure}

% Requires in preamble: \usepackage{pgfplots}, \pgfplotsset{compat=1.18}
%                       \usepgfplotslibrary{groupplots}
\begin{figure}[htbp]
\centering
\begin{tikzpicture}
\begin{groupplot}[
    group style={
        group size=2 by 1,
        horizontal sep=1.8cm,
    },
    width=0.42\textwidth,
    height=0.42\textwidth,
    grid=major,
    grid style={gray!25},
    tick align=outside,
    xlabel={Max F-bar pressure correction (ref)},
    nodes near coords,
    nodes near coords style={font=\small, anchor=south west},
    point meta=explicit symbolic,
]

%% --- Left panel: new method ---
\nextgroupplot[
    title={New method},
    ylabel={Relative error at target},
    ymax=0.7
]
\addplot[only marks, mark=*, mark size=3pt, blue] table[
    col sep=comma,
    x=fbar_pressure_correction_max_ref,
    y=rel_err_target,
    meta=nu,
]{data/strength/fbar_strength_results_new.csv};

%% --- Right panel: old method ---
\nextgroupplot[
    title={Old method},
    ylabel={Relative error at target},
    yticklabel pos=right,
    ymax=0.7
]
\addplot[only marks, mark=*, mark size=3pt, red] table[
    col sep=comma,
    x=fbar_pressure_correction_max_ref,
    y=rel_err_target,
    meta=nu,
]{data/strength/fbar_strength_results_old.csv};

\end{groupplot}
\end{tikzpicture}
\caption{Dynamic simulation - Relative error at target vs.\ maximum F-bar pressure correction for the new method (left) and old method (right), with Poisson's ratio $\nu$ labelled at each point.}
\label{fig:fbar_strength_comparison}
\end{figure}

% Requires in preamble: \usepackage{pgfplots}, \pgfplotsset{compat=1.18}
%                       \usepgfplotslibrary{groupplots}
\begin{figure}[htbp]
\centering
\begin{tikzpicture}
\begin{groupplot}[
    group style={
        group size=2 by 1,
        horizontal sep=1.8cm,
    },
    width=0.42\textwidth,
    height=0.42\textwidth,
    grid=major,
    grid style={gray!25},
    tick align=outside,
    xlabel={Max F-bar pressure correction (ref)},
    nodes near coords,
    nodes near coords style={font=\small, anchor=south west},
    point meta=explicit symbolic,
]

%% --- Left panel: new method ---
\nextgroupplot[
    title={New method},
    ylabel={Relative error at target},
    ymax=0.7
]
\addplot[only marks, mark=*, mark size=3pt, blue] table[
    col sep=comma,
    x=fbar_pressure_correction_max_ref,
    y=rel_err_target,
    meta=nu,
]{data/strength/fbar_strength_results_new.csv};

%% --- Right panel: old method ---
\nextgroupplot[
    title={Old method},
    ylabel={Relative error at target},
    yticklabel pos=right,
    ymax=0.7
]
\addplot[only marks, mark=*, mark size=3pt, red] table[
    col sep=comma,
    x=fbar_pressure_correction_max_ref,
    y=rel_err_target,
    meta=nu,
]{data/strength/fbar_strength_results_old.csv};

\end{groupplot}
\end{tikzpicture}
\caption{Dynamic simulation - Relative error at target vs.\ maximum F-bar pressure correction for the new method (left) and old method (right), with Poisson's ratio $\nu$ labelled at each point.}
\label{fig:fbar_strength_comparison}
\end{figure}

% Requires in preamble: \usepackage{pgfplots}, \pgfplotsset{compat=1.18}
%                       \usepgfplotslibrary{groupplots}
\begin{figure}[htbp]
\centering
\begin{tikzpicture}
\begin{groupplot}[
    group style={
        group size=2 by 1,
        horizontal sep=1.8cm,
    },
    width=0.42\textwidth,
    height=0.42\textwidth,
    grid=major,
    grid style={gray!25},
    tick align=outside,
    xlabel={Max F-bar pressure correction (ref)},
    nodes near coords,
    nodes near coords style={font=\small, anchor=south west},
    point meta=explicit symbolic,
]

%% --- Left panel: new method ---
\nextgroupplot[
    title={New method},
    ylabel={Relative error at target},
    ymax=0.7
]
\addplot[only marks, mark=*, mark size=3pt, blue] table[
    col sep=comma,
    x=fbar_pressure_correction_max_ref,
    y=rel_err_target,
    meta=nu,
]{data/strength/fbar_strength_results_new.csv};

%% --- Right panel: old method ---
\nextgroupplot[
    title={Old method},
    ylabel={Relative error at target},
    yticklabel pos=right,
    ymax=0.7
]
\addplot[only marks, mark=*, mark size=3pt, red] table[
    col sep=comma,
    x=fbar_pressure_correction_max_ref,
    y=rel_err_target,
    meta=nu,
]{data/strength/fbar_strength_results_old.csv};

\end{groupplot}
\end{tikzpicture}
\caption{Dynamic simulation - Relative error at target vs.\ maximum F-bar pressure correction for the new method (left) and old method (right), with Poisson's ratio $\nu$ labelled at each point.}
\label{fig:fbar_strength_comparison}
\end{figure}
